\section{树上高级算法：参数、调用与改造}

本节紧跟前面的基础模板。先看每个接口传什么，再看哪些不变量不能动，最后看如何替换维护信息或转移。

\subsection{树上启发式合并（Disjoint Set Union on Tree）}

识别信号：树不修改，每个节点有颜色/权值，每个询问问``子树中出现次数最多的颜色''\,``不同颜色数''\,``频率平方和''等。先求重儿子，保留重儿子的贡献，轻儿子算完后清空，复杂度通常 \(O(n\log n)\)。

\begin{lstlisting}[language={C++}]
vector<vector<int>> g;
vector<int> sz, son, color, ans;
vector<int> cnt;
long long cur_answer = 0; // 示例：维护所有颜色频次的平方和。

void dfs_size(int u, int fa) {
    // sz 和 son 只计算一次；son[u] 是子树最大的儿子。
    sz[u] = 1;
    for (int v : g[u]) if (v != fa) {
        dfs_size(v, u);
        sz[u] += sz[v];
        if (sz[v] > sz[son[u]]) son[u] = v;
    }
}
void add_subtree(int u, int fa, int delta) {
    // 先删除颜色旧贡献，再改变频次，最后加入新贡献。
    // 这里选 freq^2 作为演示；换题时只改这段“单点统计器”。
    cur_answer -= 1LL * cnt[color[u]] * cnt[color[u]];
    cnt[color[u]] += delta;
    cur_answer += 1LL * cnt[color[u]] * cnt[color[u]];

    // 加/删的是整棵子树，不是只有根 u。
    for (int v : g[u])
        if (v != fa) add_subtree(v, u, delta);
}
void dsu(int u, int fa, bool keep) {
    // 轻儿子先算完并清空，避免它们的统计污染当前节点。
    for (int v : g[u])
        if (v != fa && v != son[u]) dsu(v, u, false);

    // 重儿子的统计保留，之后所有轻子树都合并到这份大答案中。
    if (son[u]) dsu(son[u], u, true);

    // 把轻儿子整棵子树逐个加入。
    for (int v : g[u])
        if (v != fa && v != son[u]) add_subtree(v, u, +1);
    add_subtree(u, fa, +1);

    // 当前集合恰好是 u 的整棵子树，读取题目答案。
    ans[u] = (int)cur_answer; // 换成题目要求的统计量

    if (!keep) {
        // 父节点不需要这棵轻子树的贡献，因此整棵删除。
        add_subtree(u, fa, -1);
    }
}
\end{lstlisting}

例题：\href{https://codeforces.com/problemset/problem/600/E}{CF600E Lomsat gelral}。题目问每个子树中出现次数最多的颜色之和。上面代码中的 \passthrough{\lstinline!cur\_answer!} 要改为``维护当前最大频次 \passthrough{\lstinline!mx!}，每次 \passthrough{\lstinline!cnt[color]!} 从旧频次变为新频次时更新 \passthrough{\lstinline!sum\_color[mx]!}''；算法框架完全不变。不要把 \passthrough{\lstinline!add\_subtree!} 写成只加节点 \passthrough{\lstinline!u!}，轻儿子整棵子树必须加入。

\subsubsection{测试样例：树上启发式合并必须统计整棵子树}

为了直接验证上面示例中的 \passthrough{\lstinline!freq\^2!} 统计器，下面输出每个节点子树内各颜色频次平方和。输入每组为 \passthrough{\lstinline!n!}、颜色、\passthrough{\lstinline!n-1!} 条无向边。

\begin{lstlisting}
输入
3
1
5
3
1 2 1
1 2
1 3
5
1 2 2 3 2
1 2
1 3
3 4
3 5

输出
1
5 1 1
11 1 5 1 1
\end{lstlisting}

第二组能抓住``只加入节点 \passthrough{\lstinline!u!}、没有加入轻儿子整棵子树''的错误；第三组同时包含重儿子和轻儿子，检查清空轻子树后父节点统计是否恢复正确。


\subsection{长链剖分}

识别信号：树上大量询问``第 \(k\) 个祖先/深度为某值的节点''，\passthrough{\lstinline!k!} 很大，普通倍增内存或常数不理想。长链剖分按``向下最深''的儿子选重儿子，并把同一条长链连续存储，常配合 \passthrough{\lstinline!O(1)!} 的第 \(k\) 级祖先。

\begin{lstlisting}[language={C++}]
void dfs_len(int u, int fa) {
    // up[u] 是父亲；len[u] 是从 u 向下走的最长链节点数。
    up[u] = fa; len[u] = 1; son[u] = 0;
    for (int v : g[u]) if (v != fa) {
        dfs_len(v, u);
        // 长链剖分按“深度最长”选重儿子，不是按子树大小。
        if (len[v] + 1 > len[u])
            len[u] = len[v] + 1, son[u] = v;
    }
}
void dfs_chain(int u, int top) {
    // 同一条长链上的节点共享链顶 top，并连续编号。
    chain_top[u] = top;
    chain_pos[u] = ++timer;

    // 先沿长儿子走，保证一条长链连续。
    if (son[u]) dfs_chain(son[u], top);

    // 其他儿子都是新长链的起点。
    for (int v : g[u])
        if (v != up[u] && v != son[u]) dfs_chain(v, v);
}
\end{lstlisting}

例题：\href{https://www.luogu.com.cn/problem/P5903}{P5903 树上 K 级祖先}。读入查询 \passthrough{\lstinline!(x,k)!} 后，如果 \passthrough{\lstinline!k > depth[x]!} 输出 \passthrough{\lstinline!0!}；否则沿长链的 \passthrough{\lstinline!chain\_top!} 和预处理的链数组跳转。这个题还特别容易爆栈，\passthrough{\lstinline!n=5e5!} 时要按题目环境决定是否使用非递归 DFS 或栈扩展。

\subsubsection{测试样例：长链、星形树和超过根深度}

每组输入 \passthrough{\lstinline!n q!}，随后给出节点 \passthrough{\lstinline!2..n!} 的父亲，再给出 \passthrough{\lstinline!q!} 个 \passthrough{\lstinline!(x,k)!}；不存在的祖先输出 \passthrough{\lstinline!0!}。

\begin{lstlisting}
输入
3
1 3
1 0
1 1
1 0
4 4
1 2 3
4 0
4 1
4 3
4 4
5 4
1 1 1 1
2 1
2 2
5 0
1 0

输出
1
0
1
4
3
1
0
1
0
5
1
\end{lstlisting}

第一组检查单点和 \passthrough{\lstinline!k=0!}；第二组是最长链，第三组是星形树。\passthrough{\lstinline!k!} 超过深度时必须输出 \passthrough{\lstinline!0!}，不能访问未初始化的倍增表。


\subsection{点分治与点分树}

静态树上``距离为 \passthrough{\lstinline!k!} 的点数/距离小于等于 \passthrough{\lstinline!k!} 的点权和''是点分治；有在线修改则把每个节点到各层重心的距离放入重心树容器，形成点分树。

\begin{lstlisting}[language={C++}]
void collect(int u, int fa, int d, vector<int>& ds) {
    // ds 保存“当前重心 c 到该分支中所有节点的距离”。
    // 点分治的每层都只收集未被 removed 的节点。
    ds.push_back(d);
    for (int v : g[u])
        if (v != fa && !removed[v]) collect(v, u, d + 1, ds);
}
void decompose(int entry) {
    // get_centroid 必须只在当前未删除连通块中找重心。
    int c = get_centroid(entry); // 按子树大小找重心
    removed[c] = true;

    for (int v : g[c]) if (!removed[v]) {
        vector<int> ds;
        collect(v, c, 1, ds);

        // 把该分支的距离加入 c 的排序数组；查询时二分。
        // 同时还要把这些距离按“属于哪一个 c 的儿子分支”保存，
        // 查询时减掉同一分支的重复贡献（容斥）。
    }

    // 删除 c 后，剩下的是若干个互不相连的子问题。
    for (int v : g[c]) if (!removed[v]) decompose(v);
}
\end{lstlisting}

例题可以从\href{https://www.luogu.com.cn/problem/P6329}{P6329 点分树 \textbar{} 震波}开始：修改一个点的权值时，沿重心父链向上，把该点到每个重心的距离对应的贡献更新；查询重心时，对每层距离数组二分，并减去``同一子树方向''的重复贡献。点分树模板比普通点分治多一层``重心祖先链''，现场不要混用两个模板。

\begin{quote}
说明：当前代码只展示了点分治的分解和距离收集框架，没有给出完整的 \passthrough{\lstinline!get\_centroid!}、查询容器、修改和答案接口。因此这里不伪造``输入/标准输出''样例；等这些接口补全后，测试必须至少包含单点、链、星形树、半径为 \passthrough{\lstinline!0!} 和跨重心重复计数五类情况。
\end{quote}


\subsection{伸展树、文艺平衡树与 Link-Cut Tree 动态树}

Splay 适合维护序列：节点有左右儿子、父亲、子树大小和翻转标记；把第 \passthrough{\lstinline!k!} 个节点旋到根，就能做区间翻转、插入、删除。文艺平衡树的区间反转调用流程是：把 \passthrough{\lstinline!l-1!} 旋到根，把 \passthrough{\lstinline!r+1!} 旋到根的右儿子，目标区间就是右儿子的左子树，打 \passthrough{\lstinline!rev!} 标记。

LCT 适合动态森林的路径操作，核心接口是 \passthrough{\lstinline!makeroot(x); access(y); splay(y)!}，此时 \passthrough{\lstinline!y!} 的 splay 子树就是原树路径 \passthrough{\lstinline!x..y!}。在已有 LCT 模板中，题目操作的对应关系通常是：

\begin{lstlisting}[language={C++}]
// 路径查询 x..y：
// makeroot(x) 把 x 变成原树根；split(x,y) 再 access(y)，
// 此时 y 的辅助树包含原树路径 x..y，维护值在 tr[y] 中汇总。
makeroot(x);
split(x, y);
answer = tr[y].sum;

// 连边：先确认两点不在同一棵树，否则会形成环。
makeroot(x);
if (findroot(y) != x) tr[x].fa = y;

// 断边：把边 x-y 暴露出来后，理想形态是 y 的左儿子恰为 x，
// 且 x 没有右儿子；这两个条件都满足才能安全删除。
makeroot(x);
access(y);
splay(y);
if (ch[y][0] == x && ch[x][1] == 0)
    ch[y][0] = fa[x] = 0;
\end{lstlisting}

例题：\href{https://www.luogu.com.cn/problem/P3690}{P3690 动态树}要求路径异或、修改点权、连边和断边。先用已有板子完成 \passthrough{\lstinline!is\_root/rotate/splay/access/makeroot/findroot/split/link/cut!}，再把 \passthrough{\lstinline!sum!} 的合并从加法换成异或。LCT 第一次现场学习风险很高，遇到它时优先复制完整模板，只修改``维护信息''和输入操作映射。

\subsubsection{测试样例：动态树的路径、修改、连边和断边}

下面采用 P3690 操作编号：\passthrough{\lstinline!0 x y!} 查询路径异或，\passthrough{\lstinline!1 x y!} 连边，\passthrough{\lstinline!2 x y!} 断边，\passthrough{\lstinline!3 x v!} 修改点权。

\begin{lstlisting}
输入
3
1 3
5
0 1 1
3 1 7
0 1 1
3 7
1 2 4
1 1 2
1 2 3
0 1 3
3 2 5
0 1 3
2 2 3
0 1 2
4 7
1 1 1 1
1 1 2
1 3 4
0 1 2
1 2 3
0 1 4
2 2 3
0 1 2

输出
5
7
7
0
4
0
0
0
\end{lstlisting}

第二组检查路径合并和单点修改；第三组先建立两棵树，再连成一棵，最后断开，能抓住 \passthrough{\lstinline!makeroot/access/splay!} 顺序和断边条件错误。实际提交前还要根据题目保证避免查询不连通节点。



